\documentclass{article}
\usepackage[utf8]{inputenc}
\usepackage[T1]{fontenc}
\usepackage{geometry}
\usepackage{graphicx}
\usepackage{float}
\usepackage{hyperref}
\usepackage[english,indonesian]{babel}

\title{Optimasi Pemrosesan Template LaTeX Menggunakan Pendekatan Berbasis AI}
\author{Academic Support Agent}
\date{\today}

\begin{document}
\maketitle

\begin{abstract}

\end{abstract}

\chapter{PENDAHULUAN}


\chapter{TINJAUAN PUSTAKA}


% Original content (commented out):
% \documentclass[12pt]{article}
\usepackage[a4paper, top=2.5cm, bottom=2.5cm, left=2.5cm, right=2.5cm] {geometry} %Sayfa düzenini ayarlar.
\usepackage[utf8x]{inputenc} 
\usepackage[T1]{fontenc}
\usepackage{graphicx}
\usepackage[version=4]{mhchem}
\usepackage{caption}
\renewcommand{\figurename}{Şekil} %Şekiller için başlık düzeni
\captionsetup[figure]{skip=5pt}	%Şekil başlığından sonraki boşluk miktarı
\renewcommand{\tablename}{Tablo} %tablolar için başlık düzeni
\captionsetup[table]{skip=3pt} %tablo başlığından önceki boşluk

\renewcommand{\baselinestretch}{1.5} %Satır aralığını ayarlar.
\renewcommand{\contentsname}{İçindekiler} %Contensts yerinde içindekiler kısmını değiştirdim.

\begin{document}
\begin{titlepage}
	\newcommand{\HRule}{\rule{\linewidth}{0.5mm}} 
	
	\center %Metinleri ortalar.
			
	\textsc{\LARGE BOZOK ÜNİVERSİTESİ}\\ 
	\textsc{\Large MÜHENDİSLİK-MİMARLIK FAKÜLTESİ}\\ 
	\textsc{\Large İNŞAAT MÜHENDİSLİĞİ}\\ 
	\textsc{\Large İM426 - BETON TEKNOLOJİSİ}\\[3.5cm] 
	
	
	\HRule\\[0.4cm]
	
	{\huge\bfseries Rapor Başlığı Buraya Yazılacaktır.}\\
	
	\HRule\\[4cm]
	
	{\Large Ad Soyad: Murat Şahin}\\
	{\Large Öğrenci No: 160000000}\\
	\vfill
	
	
	{\large\today} 
\end{titlepage}


\pagenumbering{roman} %İçindekiler kısmında sayfa numaralarını roma rakamı yapar.
\tableofcontents
\newpage
\pagenumbering{arabic} %Raporun diğer kısmında sayfa numaralarını rakam yapar.

\section{Cüruflu Çimentolar} %Yeni bir bölüm açar.
Bu metin \LaTeX \  kullanım örneği olarak kullanılmak üzere Türkiye Çimento Müstahsilleri Birliği'nin sitesinden Prof. Dr. Mustafa TOKYAY ve İnş. Y. Müh. Korhan ERDOĞDU tarafından yazılmış olan  \textbf{Cüruflar ve Cüruflu Çimentolar} adlı eserden alınmıştır.

\section{Üretim}
Cüruflu çimentolar granüle yüksek fırın cürufunun portland çimentosu klinkeri
ve alçı taşı ile birlikte veya ayrı ayrı öğütülüp karıştırılmasıyla elde edilir. Cüruflu çimentonun üretimi şematik olarak Şekil 3.1’de gösterilmiştir.
\begin{figure}[h!] %Kullandığınız şeklin yerini sabitler.
  \caption{Beton Dökümü}
  \centering
  \includegraphics[width=1\textwidth]{beton.jpg} %width kullandığınız şeklin ölçeği olarak düşünebilir.
  \end{figure}\\
Cüruflu çimento üretiminde kullanılan cüruf miktarı değişik ülkelerde, çimento
tipine de bağlı olarak, önemli farklılıklar göstermektedir. Amerika Birleşik
Devletleri’nde ASTM C 595’de belirtilen Portland Yüksek Fırın Cürufu
Çimentoları (Portland Blast Furnace Slag Cement)nda cüruf oranı \%25-65’dir.
Almanya’da DIN 1164’de Eisenportland ve Hochofen çimentoları olarak
iki tip cüruflu çimento tanımlanmıştır. Bunlardan birincisinde cüruf miktarı en fazla \%40, ikincisinde ise \%41-85’dir. Fransa’da Ciment de Fer \%25-35
Ciment Métallurgique Mixte \%45-55, Ciment de Haut Foumeau \%65-75 ve
Ciment de Laitier au Clinker \%80’in üstünde granüle yüksek fırın cürufu içerir.
Ülkemizde standardı bulunan üç tür cüruflu çimentodan TS 20’de belirtilen cüruflu çimentolar CÇ 32.5 ve CÇ 42.5’da \%20-80, TS809’da belirtilen
Süper Sülfat çimentosu SSÇ 32.5’da ise en az \%65 granüle yüksek fırın
cürufu bulunur.\\[9pt]
Genel hatlarıyla, üretim enerjisi ve hammadde kullanımı değerlerine
bakıldığında, 1 ton \%65 cüruf (\%15 rutubetli) içeren bir cüruflu çimento
üretimi için 1 ton Portland Çimentosu üretiminde gerekli olan hammadde
miktarında yaklaşık \%65, kullanılan enerji miktarında ise yaklaşık \%50 bir
azalma sağlanır (1).
\section{Cüruflu Çimentoların Hidratasyonu}
Yüksek fırın cüruflarının kendi başlarına veya bazı aktivatörlerle birlikte suyla
reaksiyonlarına bölüm 2.5’de değinilmiştir. GYFClerin hidratasyonlarının
hızlandırılması için gereken iki kutup aktivatör çimentolu sistemlerde zaten
bulunmaktadır: (1) alçı ve (2) çimentonun silikat bileşenlerinin hidratasyonuyla
elde edilen Ca(OH)2. Bütün ticari cüruflu çimentolarda klinker kısmı cüruftan
daha hızlı olarak hidrate olur. Ancak, cüruf da, aşağıdaki kalorimetre
eğrilerinde de (Şekil 3.2) görüleceği gibi, çok erken yaşlarda suyla
reaksiyona girmektedir (1, 16). \%70 cüruf ve \%30 portland çimentosu
içeren bu karışımlarda 1. cürufun ikincisine göre daha reaktif olduğu
anlaşılmaktadır. Açığa çıkan hidratasyon ısılarının piklerine bakıldığında,
birinci cüruflu çimentoyla Portland çimentosunun yaklaşık aynı sürede (~6
saat) maksimum hidratasyon ısısına ulaştıkları görülür. Ayrıca, bu süre
zarfında yalnızca \%30 portland çimentosundan elde edilecek hidratasyon
ısısı 15.6 J/gh olacakken 70/30’luk cüruflu çimentolardan, birinci cüruf
için, 29 J/gh, ikinci cüruf için de 18 J/gh ısı açığa çıkmıştır. Bu durum da
cürufların daha çok erken yaşlardan itibaren hidrate olmaya başladığını
göstermektedir.\\[9pt]
Cüruflu çimentolarda cüruf ve klinker fazlarının hidratasyonuyla meydana
gelen C-S-H jelleri Tarayıcı Elektron Mikroskop ve EDAX’la incelendiğinde
cüruf taneciklerinin yüzeylerinin klinkerin hidratasyonuyla oluşan C-S-H’ın
kompozisyonuna göre \ce{AI2O3} ve \ce{MgO} bakımından daha zengin olan bir C-S-H tabakası ile kaplandığı saptanmıştır. Sülfat iyonları bu tabakayı
eriterek daha pürüzlü bir yapıya sahip ve dolayısıyla suyun penetrasyonunu
engellemeyen ikinci bir jel oluşturmasına neden olur (16, 17, 18).\\[9pt]
\ce{Ca(OH)2} aktivasyonuyla ilk olarak etringit oluşur. Daha ileri yaşlarda,
alçı tamamıyla tüketildikten sonra, etringit monosülfat formuna dönüşür ve
tetrakalsiyum aluminat hidratla (\ce{C4AH13}) katı eriyik haline gelir (1, 16).\\[9pt]
Düşük \ce{MgO} içeren cüruflu çimentolarda \ce{C3ASH14}, yüksek \ce{MgO} içerenlerde
ise \ce{C4AHX} ve bunun magnezyum analoğu olan \ce{M4AHX}, diğer hidratasyon
ürünleri olarak gözlenmiştir (19, 20, 21). Bunların dışında, cüruflu
çimentoların hidratasyonu sonucunda afwillit (\ce{C3SH3}) ve tobermorit (\ce{C3S2H3})
oluşumu da söz konusudur (22).
\begin{figure}[h!] %Kullandığınız şeklin yerini sabitler.
  \caption{Beton Dökümü}
  \centering
  \includegraphics[width=1\textwidth]{beton.jpg} %width kullandığınız şeklin ölçeği olarak düşünebilir.
  \end{figure}
\subsection{Cürufun Camsı Fazının ve Kristalinitesinin Etkileri}
GYFCnin hidratasyonuna Bölüm 2.5’de değinilmiştir. Demoulian ve
ekibine göre (6) cürufun mükemmel bir camsı yapıya sahip olması cüruflu
çimentolarda optimum reaktiviteye sahip olduklarının bir göstergesi değildir.
Bu araştırmacılar \%3-5 merwinit (\ce{C3MS2}) içeren cürufların kullanıldığı
çimentolarla en yüksek dayanımları elde etmişlerdir.
Satarin’e göre de, Bu bölüm 2.2’de sözü edilen, silika tetrahedronları cürufun
camsı fazı içinde bir tür polimer halinde bulunurlar. Eğer bu moleküllerin
polimerleşme derecesi çok yüksekse cürufun hidrolik aktivitesinde azalma
söz konusu olur (23).
Aynı portland çimentosu ve yedi değişik GYFC kullanılarak üretilen cüruflu
çimentolar üzerine yapılan bir mikro yapı araştırmasında cürufun hidratasyon
hızıyla, rekristalize hale geldikten sonra, içerdiği akermanit (\ce{C2MS2}) miktarı
arasında 0.98 korelasyon katsayısına sahip bir bağıntı saptanmıştır. Buna
göre, rekristalize cüruf içindeki akermanit miktarı ne kadar yüksekse cüruf o
kadar yavaş hidrate olmaktadır (16).
\subsection{Klinker Özellikleri ve Cüruf İnceliğinin Etkileri}
GYFCnin hidratasyonuna Bölüm 2.5’de değinilmiştir. Demoulian ve
ekibine göre (6) cürufun mükemmel bir camsı yapıya sahip olması cüruflu
çimentolarda optimum reaktiviteye sahip olduklarının bir göstergesi değildir.
Bu araştırmacılar \%3-5 merwinit (\ce{C3MS2}) içeren cürufların kullanıldığı
çimentolarla en yüksek dayanımları elde etmişlerdir.
Satarin’e göre de, Bu bölüm 2.2’de sözü edilen, silika tetrahedronları cürufun
camsı fazı içinde bir tür polimer halinde bulunurlar. Eğer bu moleküllerin
polimerleşme derecesi çok yüksekse cürufun hidrolik aktivitesinde azalma
söz konusu olur (23).
Aynı portland çimentosu ve yedi değişik GYFC kullanılarak üretilen cüruflu
çimentolar üzerine yapılan bir mikro yapı araştırmasında cürufun hidratasyon
hızıyla, rekristalize hale geldikten sonra, içerdiği akermanit (\ce{C2MS2}) miktarı
arasında 0.98 korelasyon katsayısına sahip bir bağıntı saptanmıştır. Buna
göre, rekristalize cüruf içindeki akermanit miktarı ne kadar yüksekse cüruf o
kadar yavaş hidrate olmaktadır (16).\\[9pt]
Cüruflu çimento üretiminde cüruf ve klinkerin birlikte öğütülmesi klinkerin daha
ince buna karşılık cürufun daha kaba kalmasına yol açar. Aynı öğütme
sonucunda daha ince cüruf elde etmek mümkün olur.
X-ışınları Diffraksiyonu ve Tarayıcı Elektron Mikroskop analizleri ile aynı
klinkerden üretilen bir portland çimentosunun ve \%85 cüruf içeren bir cüruflu
çimentonun hidralasyonlarının takibinden aşağıdaki sonuçlara ulaşılmıştır
(16):
1. İki günlük hidratasyon sonucunda, cüruflu çimentoda meydana gelen
etringit miktarı portland çimentosuna oranla daha fazladır. Bu fazlalğın
nedeni alçının hem klinkerdeki \ce{C3A} hem de cürufun alumina fazıyla
reaksiyona girmesi ile açıklanmıştır. Öte yandan, cüruflu çimentonun
X-ışınları difraktogramında C-S-H’a tekabül eden yaygın ve geniş pikin
altındaki alanın portland çimentosundakine göre \%30 daha fazla olması,
cürufun hidratasyonuyla da erken yaşlarda C-S-H jeli meydana geldiğinin
göstergesidir.
2. Hidratasyonun yedinci gününde etringitin kısmen monosülfat formuna
transformasyonun başladığı gözlenmiştir. Cüruflu çimentonun X-ışınları
difraktogramlarında C-S-H jeli yansımalarında artış olmuştur. Etgingit,
cüruflu çimentoda, portland çimentosuna göre, hala daha fazladır.
3. Yedi ve yirmi sekizinci günler arasında C-S-H oluşumu devam etmiştir.
Cüruf inceliğinin etkisine hızlandırılmış deneylerle bakılmış olan bir
araştırmada aynı cüruf 6310 $cm^2/g$ (Blaine) ve 3500 $cm^2/g$ inceliklerde
öğütülmüş ve 6 ve 24 saat sonunda çimentoların dayanımları belirlenmiştir
(16). Bu çalışmada portland çimentosunun inceliği sabit (3500 $cm^2/g$)
tutulmuş, cürufun aktivasyonu soda kullanılarak yapılmıştır. Deneylerde
kullanılan cüruflu çimentolar \%85 GYFC, \%13 portland çimentosu ve \%2
alçı içermektedir. Elde edilen sonuçlar Çizelge 3.1’de verilmiştir.\\
\begin{table}
\center
\begin{tabular}{|l|c|c|c|c|}\hline
Rakamlar 	& 2 & 3	& 4	& 5 \\\hline 
Rakamlar	& 2 & 3	& 4	& 5 \\\hline 
Rakamlar 	& 2 & 3	& 4	& 5 \\\hline 
Rakamlar 	& 2 & 3	& 4	& 5 \\\hline 
\end{tabular}
\caption{Deneme Tablosu}

\end{table}

\subsection{Isıl İşlemlerin Etkileri}
Bütün kimyasal reaksiyonlarda olduğu gibi, ortam sıcaklığının artması
hidratasyon hızının artmasına neden olur. Hidratasyon sırasında meydana
gelen reaksiyonlar için gerekli alan aktivasyon enerjisi cüruflu çimentolarda,
eşdeğer portland çimentolarına göre yaklaşık \%8 daha yüksektir (16, 24,
25). Bu nedenle, cüruflu çimento içeren sistemlere ısıl işlem uygulanması,
dışarıdan ilave bir enerji kaynağı sağlayarak hidratasyonu kolaylaştırıcı ve
hızlandırıcı bir yöntem olarak kabul edilmektedir.
Ayrıca, çok yüksek sıcaklıklar ve basınç altında otoklavlama yöntemi
uygulanarak yaptırılan hidratasyon cüruflu çimentoların dayanımlarını çok
artırmaktadır (16).
\section{Fiziksel Özellikler}
Portland Yüksek Cüruflu Çimentoları Portland Çimentosundan daha açık
renklidir. Özgül ağırlığı, Portland Çimentolarına oranla biraz düşüktür (3.00-
3.05). Priz süreleri, rötre, sünme, betonarmede donatı aderansı özellikleri
bakımından portland yüksek fırın cürufu çimentolarıyla portland çimentoları
arasında önemli farklılıklar yoktur (4).
Portland Yüksek Fırın Cürufu çimentolarının hidratasyon ısıları kullanılan
klinkerin ve GYFC nin kompozisyonu ve miktarıyla orantılı olarak değişir
ancak, genel olarak, normal portland çimentolarınınkinden daha düşüktür
(4,7).
\section{Dayanım Özellikleri}
Cüruflu çimentolar hidratasyon mekanizması ve hidratasyon ürünleri
bakımından portland çimentolarıyla benzer nitelikler göstermekle birlikte,
normal incelikte öğütülmüş durumda dayanım kazanma hızı, özellikle erken
yaşlarda, portland çimentolarından daha düşüktür. Aynı 28 günlük dayanımın
elde edildiği, birbirine çok yakın inceliklerdeki (3000 $cm^2/g$) bir portland
çimentosuyla \%65 cüruf içeren bir cüruflu çimento karşılaştırıldığında,
erken yaşlarda portland çimentosunun geç yaşlarda ise cüruflu çimentonun
daha yüksek dayanıma sahip olduğu bilinmektedir (1,5). Durum betonda
incelendiğinde, eğilme dayanımlarının \%25-45 cüruf içeren çimentolarla
yapılan betonlarda, normal portland çimentosuyla yapılanlara göre, ya eşit
ya daha yüksek olduğu saptanmıştır (26).
Cüruflu çimentoların uygun erken dayanımlara sahip olabilmesi için özgül
yüzeylerinin ve cüruf miktarlarının ayarlanması gerekir. Erken yaşlarda
dayanımın yükseltilmesi için uygulanan bir başka yöntem de ağırlıkça \%0.5
miktarında trietanolamin kullanılmasıdır (27).
Cüruf inceliğinin etkilerine bölüm 3.2.2’de değinilmiştir. Cürufun klinkere
oranla daha zor öğütülebilir bir malzeme olması, birlikte öğütme koşullarında,
klinkerin ince cürufun ise daha kaba kalması durumunu ortaya çıkarmaktadır.
Ayrı öğütme koşullarında elde edilmiş, \%30-50 cüruf içeren çimentolarla
yapılan betonlarda gerek işlenebilirlik gerekse erişen dayanım bakımından
önemli artışlar sağlanmıştır (28).
Cüruflu çimentoların dayanımlarının tahmini için önerilen çeşitli Hidrolik
İndeksler Bölüm 2.4’de verilmiştir. Bunların yanısıra, iki Japon araştırmacı
tarafından önerilen bir çoklu regresyon denkleminde cüruflu çimentoyla aynı
klinkerden elde edilmiş olan bir normal portland çimentosunun 28 günlük
dayanımlarının oranı (Y) cürufun alkalinitesine (X1), camsı faz miktarına (X2)Cüruflu çimentolar hidratasyon mekanizması ve hidratasyon ürünleri
bakımından portland çimentolarıyla benzer nitelikler göstermekle birlikte,
normal incelikte öğütülmüş durumda dayanım kazanma hızı, özellikle erken
yaşlarda, portland çimentolarından daha düşüktür. Aynı 28 günlük dayanımın
elde edildiği, birbirine çok yakın inceliklerdeki (3000 $cm^2/g$) bir portland
çimentosuyla \%65 cüruf içeren bir cüruflu çimento karşılaştırıldığında,
erken yaşlarda portland çimentosunun geç yaşlarda ise cüruflu çimentonun
daha yüksek dayanıma sahip olduğu bilinmektedir (1,5). Durum betonda
incelendiğinde, eğilme dayanımlarının \%25-45 cüruf içeren çimentolarla
yapılan betonlarda, normal portland çimentosuyla yapılanlara göre, ya eşit
ya daha yüksek olduğu saptanmıştır (26).
Cüruflu çimentoların uygun erken dayanımlara sahip olabilmesi için özgül
yüzeylerinin ve cüruf miktarlarının ayarlanması gerekir. Erken yaşlarda
dayanımın yükseltilmesi için uygulanan bir başka yöntem de ağırlıkça \%0.5
miktarında trietanolamin kullanılmasıdır (27).
Cüruf inceliğinin etkilerine bölüm 3.2.2’de değinilmiştir. Cürufun klinkere
oranla daha zor öğütülebilir bir malzeme olması, birlikte öğütme koşullarında,
klinkerin ince cürufun ise daha kaba kalması durumunu ortaya çıkarmaktadır.
Ayrı öğütme koşullarında elde edilmiş, \%30-50 cüruf içeren çimentolarla
yapılan betonlarda gerek işlenebilirlik gerekse erişen dayanım bakımından
önemli artışlar sağlanmıştır (28).
Cüruflu çimentoların dayanımlarının tahmini için önerilen çeşitli Hidrolik
İndeksler Bölüm 2.4’de verilmiştir. Bunların yanısıra, iki Japon araştırmacı
tarafından önerilen bir çoklu regresyon denkleminde cüruflu çimentoyla aynı
klinkerden elde edilmiş olan bir normal portland çimentosunun 28 günlük
dayanımlarının oranı (Y) cürufun alkalinitesine ($X_1$), camsı faz miktarına ($X_2$)ve Cüruflu çimentonun pikrik asit-metanol karışımında eritildikten sonraki
kalıntı miktarına ($X_3$) bağlı olarak verilmiştir (29):
\begin{equation}
Y= -73,33 + 74,06X_1 + 0,482X_2 - 2,41X_3
\end{equation}

\section{Dayanıklılık Özellikleri}
Sülfatlı sular, deniz suları, klorlu sular, karbonatlı sular, termal sular, buz
çözücü maddeler vb. ile yapılan uzun süreli deneyler sonucunda cüruflu
çimentoların zararlı kimyasal etkiler altında performanslarının yüksek olduğu
belirlenmiştir (1).
\subsection{Sülfat Etkisine Dayanıklılık}
Regourd (16) tarafından yapılmış olan bir çalışmada sentetik olarak
hazırlanmış, \%50 \ce{C3S}+\%15 \ce{C3A}+\%5 alçı+\%30 kuartz içeren numuneler %5
\ce{MgSO4} eriyiğine batırılmıştır. İlk günden başlayarak, numunelerde genleşme
gözlenmiş ve yedinci günde tüm numuneler dağılmıştır. daha sonra bu
numuneler X-Işınları Difraksiyonu ve Tarayıcı Elektron Mikroskop kullanılarak
incelendiğinde, \ce{Mg(OH)2}, \ce{CaSO4.2H2O}, etringit ve C-M-S-H oluştuğu
saptanmıştır. Aynı araştırmada, ikinci seri numuneler kuartz yerine GYFC
kullanılarak hazırlanmış ve aynı koşullara tabi tutulmuştur. Sonuçta, GYFC
içeren numunelerde hiçbir bozulma görülmemiştir. X-Işınları Difraksiyonu ve
Tarayıcı Elektron mikroskop analizlari bu numunelerde yoğun bir CSH, bir
miktar Ca(OH)2, monosülfat formu ve çok az etringit oluştuğunu göstermiştir.
Benzer sonuçlar \%60 cüruf içeren çimentolarla yapılan beton numuneler
üzerinde de elde edilmiştir (30). Aynı araştırmadan çıkan bir başka sonuç ise
kullanılan cürufun inceliğinin artmasıyla, betonun porozitesini azalttığından
dolayı, kimyasal etkilere karşı direncin yükselmesidir.
\ce{MgSO4}, \ce{CaSO4} ve \ce{Na2SO4} eriyiklerinin kullanıldığı benzeri bir başka
araştırmada da 130 değişik çimento kullanılmış ve yüksek fırın cüruflarının,
miktarlarıyla orantılı olarak, sülfat direncini artırdığı saptanmıştır (31).
\subsection{Klor Etkisine Dayanıklıklık}
Cüruflu çimentoların klor direnci konusunda çok sayıda değişik deney
yöntemlerinin uygulandığı araştırma yapılmıştır. Bunlardan bazıları burada
özetlenmiştir.
Hızlandırılmış Klor Permeabilite deneyi uygulanarak yapılan bir araştırmada
değişik inceliklerde ve değişik cüruf miktarı içeren çimentolar kullanılarak
102mm çap ve 52mm boyda silindir harç numuneleri hazırlanmış ve
14 günlük standard bakımdan sonra numunelerin bir yüzü sodyum klorür
diğer yüzü sodyum hidroksit eriyiklerine batırılarak 60 V sabit potansiyel
farkı uygulanmış ve bir yüzden öbürüne geçen elektriksel akım ölçülerek
klor permeabilitesiyle ilişkilendirilmiştir. Bu araştırmalarda kullanılan kontrol
çimentosu (portland çimentosu) 0.865 $m^2/g$ (BET) kaba cüruflu çimentolar
0.866 $m^2/g$ (BET), orta cüruflu çimentolar 1.083 $m^2/g$ (BET) ve ince cüruflu
çimentolar 1.337 $m^2/g$ (BET) inceliktedir. İnce cüruflu çimentolarda ağırlıkça
\%70 ve \%85, orta cüruflu çimentoda \%50 ve kaba cüruflu çimentolarda
\%50, \%70 ve\%85 cüruf kullanılmıştır. elde edilen sonuçlar grafik olarak
Şekil 3.3’de gösterilmiştir (32).\\[9pt]
Benzer bir araştırma da \%8 silis dumanı, \%25 uçucu kül, \%50 GYFC içeren
numuneler üzerinde yapılmış ve cürufun diğer mineral katkılara göre daha
etkin olduğu saptanmıştır (33).
Klor permeabilitesinin tayini ile ilgili ilginç bir başka yöntemde ise
100x200x300mm boyutlarındaki 28 günlük beton plakaların üst yüzeyine
\%5 NaCI eriyiği göllendirme suretiyle 4 gün uygulanmakta, bunu takip eden
3 günde ise saf su uygulanmaktadır. Söz konusu çevrim 48 hafta sürdürülmekte
ve daha sonra yüzeyden 10, 20, 30, 40 ve 50 mm derinliklerden numuneler
alınıp klor analizi yapılmaktadır. Bu deney yönteminin uygulandığı bir
araştırmada bir sülfata dayanıklı çimento, bir normal portland çimentosu ve
\%40, 50 ve 70 GYFC içeren üç cüruflu çimento kullanılmıştır. Bu araştırmanın
sonucunda yüksek fırın cürufu kullanımının klor permeabilitesini azalttığı ve
kullanılan cüruf miktarı arttıkça bu etkinin daha belirginleştiği saptanmıştır.
Elde edilen sonuçlar 3.4’de gösterilmiştir (34).
\subsection{Deniz Suyu Etkisine Karşı Dayanıklılık}
Cüruflu çimentoların deniz suyu etkisine maruz betonlardaki yüksek
performansı yıllardır bilinmektedir. Deniz suyuna tamamıyla batırılmış
40x40x160mm harç numunelerle yapılan deneylerde kullanılan cüruf
miktarının etkisi araştırılmış ve çimento içindeki GYFC miktarının artmasıyla
genleşmenin azaldığı saptanmıştır (1, 16). Bu araştırmadan elde edilen
sonuçlar Çizelge 3.2 ve 3.3’de özetlenmiştir.\\[9pt]
Diğer bir araştırmada ise 100x50x50 mm beton numuneler yarı boylarına
kadar deniz suyuna batırılarak deneyler yapılmıştır (31). Bu araştırmada
ASTM Tip V’ye tekabül eden 24, Tip II’ye tekabül eden 16, Tip I’e tekabül
eden 17, \%15 GYFC, uçucu kül ve tras ieren 5’er \%30 cüruf içeren 12
ve \%70’in üstünde cüruf içeren 35 çimento kullanılmıştır. Yapılan deneyler
sonucunda en yüksek performans \%70’in üstünde cüruf içeren çimentolarla
elde edilmiştir.
\subsection{Karbonatlaşmaya Dayanıklılık}
Betonda karbonatlaşma açısından, beton yeterince yoğun ve ortamda yeterli
rutubet var olduğu sürece, portland çimentosuyla cüruflu çimento arasında
bir fark yoktur (1, 5, 16). Longuet (35) tarafından yapılan bir araştırmada
\%80’in üstünde cüruf içeren çimentoların, karbonatlaşmanın neden olduğu
deposivasyon sonucunda betonarme donatıların paslanmasını portland
çimentoları kadar etkin bir şekilde önlediği saptanmıştır. Öte yandan, bir
başka araştırmada, daha kısa sürelerle bakımı yapılan (1, 3 ve 7 gün) harç numunelerde yukarıda belirtilenen tam tersi sonuçlar alınmıştır. Sözü edilen bu araştırmada, karbonatlaşma direncinin artırılabilmesi için bakım süresinin uzatılması gerektiği belirtilmiştir (36).
\subsection{Alkali-Agrega Reaksiyonu}
Betonda karbonatlaşma açısından, beton yeterince yoğun ve ortamda yeterli
rutubet var olduğu sürece, portland çimentosuyla cüruflu çimento arasında
bir fark yoktur (1, 5, 16). Longuet (35) tarafından yapılan bir araştırmada
\%80’in üstünde cüruf içeren çimentoların, karbonatlaşmanın neden olduğu
deposivasyon sonucunda betonarme donatıların paslanmasını portland
çimentoları kadar etkin bir şekilde önlediği saptanmıştır. Öte yandan, bir
başka araştırmada, daha kısa sürelerle bakımı yapılan (1, 3 ve 7 gün) harç
numunelerde yukarıda belirtilenen tam tersi sonuçlar alınmıştır. Sözü edilen
bu araştırmada, karbonatlaşma direncinin artırılabilmesi için bakım süresinin
uzatılması gerektiği belirtilmiştir (36).\\[9pt]
Cüruflu çimentolar kullanılarak yapılan çok sayıda araştırma cüruf miktarının
artmasıyla alkali-agrega reaksiyonunun neden olduğu genleşmelerin
azaldığını göstermiştir (1, 16, 37). Yüksek miktarda alkali içeren bir portland
çimentosunu \%24 ve \%50 oranlarıyla GYFC ile ikame edilmesi sonucunda
ulaşılan alkali-agrega reaksiyonu genleşmeleri karşılaştırmalı olarak Şekil
3.5’de verilmiştir (1).
\subsection{Donma Çözülme Direnci}
Yüksek Fırın Cüruflu çimentoların betonların donma-çözülme direncine
etkileri, dayanım ve betonun hava miktarı sabit tutulduğu sürece, portland
çimentolarınınkinden farklı değildir. Ancak, cüruf miktarı çok yüksek olduğu
taktirde, az bir miktar düşüş görülebilir (38).\\[9pt]
Öte yandan, yüksek fırın cürufu kullanımının beton içindeki gözenek
boyutlarında, gerek fiziksel gerekse hidratasyon sonucunda, azalmaya
neden olması betonun donma çözülme direncini yükselttiği görüşü de öne
sürülmektedir (39).\\[9pt]
Yapılan tüm araştırmalarda hemfikir olunan bir sonuç cüruflu çimentoların,
betonda sabit bir hava miktarı sağlamak için, portland çimentolarına göre
daha fazla hava sürükleyici katkıya ihtiyaç gösterdikleri hususudur (1, 38,
39, 40).
\newpage %KAynakların yeni sayfada olmasını sağlar.
\section*{Kaynaklar}
\addcontentsline{toc}{section}{\protect\numberline{}Kaynaklar}%Kaynakları İçindekilere eklemek için kullanıldı.
\begin{enumerate}
\item Tokyay, M. (2016). \textit{Cement and concrete mineral admixtures}. Boca Raton: CRC Press, Taylor \& Francis Group.
Straub, C., Quercia, G., Florea, M. V. A., Brouwers, \item Straub, C., Quercia, G., Florea, M. V. A., Brouwers, H. J. H. (2014). Slump flow of Autoclaved Aerated Concrete slurries. \textit{International Conference on Non-Tradional Cement and Concrete}, Brno, Czech Republic.
\item Rashad, A. M., Bai, Y., Basheer, P. A. M., Milestone, N. B., Collier, N. C. (2013). Hydration and properties of sodium sulfate activated slag. \textit{Cement and Concrete Composites, 37}, 20-29.
\end{enumerate}
\end{document}


\end{document}